% ============================================================
%  ch02a：§2 引言、2.1 算术与几何朗兰兹（原文 pp.5–9）
% ============================================================
\section{AGKRRV 理论}

我们首先概述 [2] 与 [4] 的工作以及一些背景材料。

这些工作诚然是技术性的。我们不打算在这里概述这些工作所需的每一项技术背景。不过，我们尽量解释某些技术问题为何出现（例如，为何需要导出代数几何）。但在这方面，我们自由地诉诸本主题中一些基础性的想法，它们可能并非所有读者都熟悉：叠、DG 范畴、$D$-模、$\ell$-进层与 $\IndCoh$ 都属此类。尽管这些主题技术性强且并不总是广为人知，但如今已有许多参考文献（以及乐于助人的专家），我们认为有兴趣的读者应当不难找到资源，以探究讨论中出现的背景材料。

本节结构如下。首先，在 §2.1 中，我们稍微解释对自守函数感兴趣的读者应如何看待 de Rham（或 $D$-模）几何朗兰兹猜想，并强调后者相对于前者的若干可喜特征。在 §2.3 中，我们解释受限几何朗兰兹对应；其中一侧涉及 $\ell$-进层的模理论，我们在 §2.2 中解释。在 §2.4 中，我们解释当考虑 Frobenius 时，在有限域上故事如何展开。最后，在 §2.5 中，我们描述我们的主要算术结果（来自 [4]）如何被证明，并强调 Xue 关于 shtuka 上同调层的工作所起的关键作用。

\subsection{算术与几何朗兰兹}

\pnumt{算术朗兰兹。}{} 传统的算术朗兰兹关乎\emph{自守表示}，按定义，它们是出现在某个合适的自守函数空间中的某些\emph{不可约}表示。

朗兰兹哲学的一种粗糙（也许有些俗气）形式预言：$G$ 的自守表示对应于 $G$ 的朗兰兹参数。这里有一些修正，但它们并非我们的重点：应当使用 Arthur 参数，会出现 $L$-包，对一般数域并没有朗兰兹参数（或朗兰兹群）的合适定义，等等。我们的重点是该猜想的\emph{原子性}：基本对象是某个函数空间的不可约子商，而不是函数空间本身。

\pnumt{几何朗兰兹。}{} 与此相对，\emph{几何}朗兰兹猜想的传统形式预言：
\begin{equation}
  \Dmod(\Bun_G) \simeq \IndCoh_{\Nilp^{\mathrm{spec}}}(\LS^{\dR}_{\bG}) \simeq \QCoh(\LS^{\dR}_{\bG}).
\end{equation}
在上式中，$\IndCoh_{\Nilp^{\mathrm{spec}}}$ 是 $\QCoh$ 的一个合适扩张，由 Arinkin--Gaitsgory 定义，我们将在 §2.1.11 中稍作讨论。

由于本次会议主要关心朗兰兹纲领的\emph{算术}方面，我们花一些时间解释\emph{几何}设定的若干起始特征，包括上面使用的记号以及思考那里出现的主要对象的一些方式。

几何朗兰兹猜想的这种形式归功于 Beilinson--Drinfeld 与 Arinkin--Gaitsgory，参见 [1] 与 [13] 中对这一思想圈的导引。我们将把它称为几何朗兰兹猜想的 \emph{de Rham} 形式，因为 $D$-模理论记忆的是 de Rham 上同调群。

\pnum{}几何朗兰兹猜想的输入是一条固定域 $k$ 上的光滑射影曲线 $X$。为简化某些要点，我们假定 $k$ 代数闭，尽管在 de Rham 设定中这并非本质性的。

于是 $\Bun_G$ 是 $X$ 上 $G$-丛的空间。更确切地说，$\Bun_G = \Bun_G(X)$ 是一个\emph{叠}，其点函子由
\[
  \Bun_G(S) := \Hom(X \times S, \mathbf{B}G) = \{X \times S \text{ 上的 } G\text{-丛}\}
\]
给出，其中 $S$ 是仿射概形，$\mathbf{B}G$ 是 $G$ 的分类叠。众所周知，$\Bun_G$ 是局部有限型的光滑代数叠，尽管它不是拟紧的。

Weil 的一个漂亮公式\footnote{这个公式基本上是显然的，只要知道 $G$-丛在光滑射影曲线上是 Zariski（而不只是 étale）局部平凡的。对 $GL_n$，这由下降得到。对其他群，这是 Steinberg 定理。我们在此指出，该定理在有限域上也成立，这一点在主题中常被默认。对单连通 $G$，这是 Harder [19] 的定理。一般地，取满射 $G'\to G$，其中 $G'$ 的导群单连通、$\operatorname{Ker}(G'\to G)$ 是（连通）环面；于是 $X$ 上的 $G$-丛可提升到 $G'$（由类域论的 Brauer 群考虑），从而化归为 Harder 定理。} 说：
\begin{equation}
  \Bun_G(k) = G(F)\backslash G(\A)/G(\OO)
\end{equation}
其中 $\OO = \prod_{x \in X(k)} \OO_x$ 是\emph{整 adele} 环（$\OO_x$ 是以 $x \in X(k)$ 为基点的 Taylor 级数环），$\A = \colim_{S \subseteq X(k)\ \text{有限}} \bigl(\prod_{x \in S} \mathrm{Frac}(\OO_x) \times \prod_{x \notin S} \OO_x\bigr)$ 是类似定义的 adele 环，而 $F = k(X)$ 是 $X$ 上的有理函数域。因此，我们可以把 $\Bun_G$ 看作双重商空间的一个几何化身，无分歧自守函数就生活在那里（如果把 $k$ 换成有限域的话）。
\pnum{}几何表示论中的一个基础类比说：当 $k$ 特征为零时，叠 $\mathcal{Y}$ 上的 $D$-模范畴的行为类似于集合 $\mathcal{Y}(\F_q)$ 上的函数空间——如果这样的东西有意义的话。\footnote{例如，读者可以设想 $\mathcal{Y}$ 定义在 $\ZZ[1/N]$ 上，$N$ 与 $q$ 互素。但我要劝读者不要在这个点上过于字面化。}

这个想法有几个理由。首先，对 $k = \mathbb{C}$，某些 $D$-模通过 Riemann--Hilbert 对应与构造可构造层相关，而后者又通过 Riemann 存在定理与 \et{} 层相关；对 $k = \F_q$，后者又通过 Grothendieck--Deligne 的层—函数对应与函数相关。\footnote{我们参考 [28] 的前几节“Applications de la formule des traces aux sommes trigonométriques”以了解这个概念的背景。}

或者，人们可以设想，$D$-模编码线性微分方程组，其解在 $\mathcal{Y}(\mathbb{C})$ 上定义函数（或分布），而这些函数以不同理由类似于 $\mathcal{Y}(\F_q)$ 上的函数。

在实践中，在这个类比里使用 $\mathcal{Y}$ 上的\emph{全部} $D$-模很重要。例如，此设定中的 \emph{Mellin 变换}是一个等价
\[
  \Dmod(\mathbb{G}_m) \simeq \QCoh(\mathbb{A}^1/\ZZ);
\]
它可以看作几何朗兰兹型等式的简化玩具模型。在 Mellin 变换下，左端的全纯对象与正则全纯对象在右端都没有合理的描述。人们把这视为一个信号：在几何表示论中应当使用“全部”$D$-模的范畴，而不是某种可构造的子范畴。

此外，由 [24]，对 $G = \mathbb{G}_m$，等价 (2.1) 并不来自 Abel 范畴的等价；也就是说，在这个类比中必须使用\emph{导出}范畴。按照现代理解，我们使用 $\infty$-范畴的同伦形式中的 DG 范畴；我们一般遵守这样的约定：我们的 DG 范畴应当具有直和，而它们之间的函子应当是线性的、正合的，并保持直和。同伦形式的好处在于，它通过引入代数工具\emph{减轻}了主题的基础负担——在这种语言中，我们可以最自如地谈论幺半范畴、模范畴、张量积等等。

\begin{example}
按照前面的讨论，人们把 $\Dmod(\Bun_G)$ 看作无分歧自守函数空间 $\Aut^{\unr}_{\mathrm{c}}$ 的类比物。
\end{example}

\pnum{}我们再停留片刻，进一步讨论范畴与向量空间之间的类比。

其起源可以这样理解：对定义在 $\F_q$ 上的叠 $\mathcal{Y}/\F_q$，一个可构造 Weil 层 $\mathcal{F}$ 通过取 Frobenius 在有理点纤维上的迹，给出 $\mathcal{Y}(\F_q)$ 上的一个函数，这给出了从层产生函数的一个相当一般的程序。这就是层与函数之间类比的通常来源。

人们可以说，空间上的函数构成一个向量空间，而空间上的层构成一个范畴，于是（函数的）向量空间被范畴化（为层的范畴）。

\pnum{}我们可以把前面的讨论变得更精确。

从根本上说，前面讨论中函数的来源是：如果有一个（有限维）向量空间 $V$ 连同线性变换 $T: V \to V$，我们可以构造 $\mathrm{tr}_V(T)$ 而得到一个数。

类似地，对一个（可对偶化的 DG）范畴 $\mathcal{C}$ 连同自函子 $T: \mathcal{C} \to \mathcal{C}$，存在一个与这份数据相伴的迹 $\mathrm{tr}_{\mathcal{C}}(T) \in \Vect$；关于这个构造的详细讨论，我们参考 [15]。我们想强调：除了围绕范畴的（也许是严重的）心理障碍之外，这个一般构造相当形式化，并且与通常的迹理论相镜像。

\pnum{}我们现在转而更仔细地解释几何朗兰兹等价。

空间 $\LS^{\dR}_{\bG} = \LS^{\dR}_{\bG}(X)$ 是 $X$ 上 de Rham $\bG$-局部系统的模叠。在本文中，叠 $\LS^{\dR}_{H}$（对 $H/k$ 仿射代数群）按惯例定义为具有如下 $S$-点：
\[
  \LS^{\dR}_{H}(S) := \Hom(X_{\dR} \times S, \mathbf{B}H),
\]
其中 $X_{\dR}$ 是 $X$ 的 de Rham 空间。在这里详细讨论 de Rham 空间过于冗长，但要点是 $\QCoh(X_{\dR}) = \Dmod(X)$。

不那么循惯例地，也可以如下进行。首先，在 $k$-点的层面上：一个 $H$-局部系统应当是什么？我们可以取 $\qLisse^{\dR}(X) \subseteq \Dmod(X)$ 为其每个上同调都是局部系统（即带联络的向量丛）之余极限的对象组成的子范畴；这是 \emph{lisse} $D$-模的一个合适导出范畴，但我们按 [2] 的惯例称它们为\emph{拟 lisse}。

于是 de Rham $H$-局部系统本质上是一个对称幺半函子 $\Rep(H) \to \qLisse^{\dR}(X)$；对 $H = \mathbb{G}_m$ 这还不太对，因为此时这样一份数据是 $\qLisse^{\dR}(X)$ 中的一个张量可逆对象 $\sigma$，即一个\emph{上同调平移过的}带联络线丛；为消除这种歧义，我们细化定义，要求我们的函子是（在作一次上同调平移 $[1]$ 之后）右 $t$-正合的（以照顾到该函子把平凡表示送到常层 $\mathbf{e}_X$ 这一事实，后者位于次数 $\dim(X) = 1$）。\footnote{当然，这个问题之所以出现，只是因为我们坚持使用导出范畴——读者可以对此提出异议。在 §2.1.12 中，我们解释在我们所讲的故事中必须使用导出代数几何，因此我们的仿射测试概形 $S$ 也应当是导出的；导出范畴 $\QCoh(S)$ 对导出几何敏感，而 Abel 范畴 $\QCoh(S)^\heartsuit$ 则不然。

鉴于这些问题，我们作了一个教学上的选择：坚持使用导出范畴与（相差平移的）右 $t$-正合函子。}\footnote{读者可能会问：右 $t$-正合函子 $\Rep(H) \to \qLisse^{\dR}(X)$ 也是 $t$-正合的（两种情形都相差平移），那为什么还要写“右”？原因是，我们很快会在 (2.3) 中推广这一设定。譬如说那里的 $H$ 是平凡群，于是我们谈论的是对称幺半函子 $\Vect \to \QCoh(S)$，它把 $k$ 送到 $\mathbf{e}_X \boxtimes \OO_S$，其中 $\mathbf{e}_X$ 是常层。在平移 $\dim(X)$ 之后，后者总是连通的（即位于上同调次数 $\leq 0$），但只有当 $S$ 是经典概形（即“通常”代数几何的对象，而非导出代数几何的对象）时，它才位于 $t$-结构的核心中。}

我们注意到，这个定义的行为本质上符合预期：一个 de Rham $GL_n$-局部系统是 $X$ 上秩 $n$ 的带联络向量丛；一个 de Rham $SO_n$-局部系统是秩 $n$ 的带联络向量丛 $\mathcal{E}$ 连同保持联络的非退化对称配对 $\mathcal{E} \otimes_{\OO_X} \mathcal{E} \to \OO_X$；de Rham $Sp_{2n}$-局部系统类似，但非退化配对是反对称的；de Rham $G_2$-局部系统是带联络的 octonion 丛；等等。

这种局部系统的 Tannaka 定义——这也许是对一般代数群定义局部系统的最简单方式——经适当调整可给出 $\LS^{\dR}_{H}$ 的 $S$-点：即对称幺半函子
\begin{equation}
  \Rep(H) \to \Dmod(X) \otimes \QCoh(S),
\end{equation}
它们在相差 $\dim(X) = 1$ 的平移的意义下是右 $t$-正合的。

我们建议读者参考 [2] §4.1 以了解进一步的相关讨论。我们简要指出，任何对称幺半函子 $\Rep(H) \to \Dmod(X)$ 都落在 $\qLisse^{\dR}(X)$ 中，尽管 (2.3) 一般不会映入 $\qLisse^{\dR}(X) \otimes \QCoh(S)$。

\pnum{}不必多说：对数论学家而言，$\LS^{\dR}_{\bG}$ 被看作 Galois 参数的模叠。

确实，由 Riemann--Hilbert 对应（的一个很容易的形式），对 $k = \mathbb{C}$，存在 $\bG$-局部系统与同态 $\rho: \pi_1(X(\mathbb{C})) \to \bG(\mathbb{C})$（模共轭）之间的解析等同（尽管这在 $S$-族中不能天真地成立）。

\pnum{}几何朗兰兹猜想中下一块关键结构是\emph{谱作用}。

这是幺半范畴 $\QCoh(\LS^{\dR}_{\bG})$ 在 $\Dmod(\Bun_G)$ 上的作用，由 Drinfeld--Gaitsgory 构造。根据\emph{同上}，这个作用由其与 $\Dmod(\Bun_G)$ 上 Hecke 作用（的一个足够强的版本）的相容性唯一刻画。

在此我们参考 [?] §1.1 中关于唯一性的讨论，以及 [?] §1.5 中与 Hecke 函子的确切联系。% [Tutor Note: 原文两处均作 [?]，据上下文应为 [2]。]

\pnum{}谱作用可以如下可视化。

令 $\mathcal{Y}$ 为任意叠，并设 $\mathcal{C}$ 是 $\QCoh(\mathcal{Y})$ 的一个模范畴。我们把它画成 $\mathcal{Y}$ 上的一个纤维化范畴：
\begin{equation*}
\begin{tikzcd}[column sep=large, row sep=large]
  \mathcal{C}_y \arrow[r] \arrow[d] & \mathcal{C} \arrow[d] \\
  y \arrow[r, hook] & \mathcal{Y}
\end{tikzcd}
\end{equation*}
这里 $y$ 处的纤维 $\mathcal{C}_y$ 定义为：
\[
  \mathcal{C}_y := \mathcal{C} \underset{\QCoh(\mathcal{Y})}{\otimes} \Vect,
\]
如果 $y$ 是 $k$-点；如果它是 $A$-点，则把 $\Vect$ 换成 $A$-$\mathrm{mod}$。启发式地，我们可以写 $\mathcal{C} = \int_{y \in \mathcal{Y}} \mathcal{C}_y \, dy$。这个形式体系在 [14] 中有详细研究。

于是，非形式地说，Drinfeld--Gaitsgory 谱作用说的是：范畴 $\Dmod(\Bun_G)$ 纤维化于 $\LS^{\dR}_{\bG}$ 之上，并且这一结构由 Hecke 函子典范地决定。因此，谱作用的\emph{存在性}可以解释为自守层范畴（这个说法在此 de Rham 语境中指的是 $\Dmod(\Bun_G)$）的（范畴式）\emph{互反律}。

按定义，$\Dmod(\Bun_G)$ 在点 $\sigma \in \LS^{\dR}_{\bG}$ 处的纤维 $\Dmod(\Bun_G)_\sigma$ 是特征值为 $\sigma$ 的 \emph{Hecke 特征层}范畴。

\pnum{}在上面启发式的公式 $\mathcal{C} = \int_{y \in \mathcal{Y}} \mathcal{C}_y \, dy$ 中，我们设想在 $\mathcal{Y}$ 上有一个范畴值测度 $\mathcal{C}_y \, dy$。在我们关心的情况中，我们可能希望计算它。

这正是完整几何朗兰兹猜想的工作。$\QCoh(\LS^{\dR}_{\bG})$-模范畴 $\IndCoh_{\Nilp^{\mathrm{spec}}}(\LS^{\dR}_{\bG})$ 在这个隐喻中编码 Plancherel 测度的类比物。

这里，具有幂零奇异支撑的\emph{非相干}层范畴 $\IndCoh_{\Nilp^{\mathrm{spec}}}(\LS^{\dR}_{\bG})$ 在 [1] 中被定义，此后一直是该领域广泛研究的对象。它是 $\QCoh$ 的一个修改，具有反映 $\LS^{\dR}_{\bG}$ 的奇点某些信息的几何性质。关于这个主题的导引，我们参考 [1]。我们用 $\Nilp^{\mathrm{spec}} \subseteq T^*[-1]\LS^{\dR}_{\bG}$ 表示\emph{谱幂零锥}，并提醒它在许多其他文献中常简记为 $\Nilp$。

由于不可约 $\bG$-局部系统不支持其伴随丛的非零幂零水平截影，我们有：
\[
  \QCoh(\LS^{\dR,\mathrm{irred}}_{\bG}) \underset{\QCoh(\LS^{\dR}_{\bG})}{\otimes} \IndCoh_{\Nilp^{\mathrm{spec}}}(\LS^{\dR}_{\bG}) = \QCoh(\LS^{\dR,\mathrm{irred}}_{\bG})
\]
（详见 [1] 命题 13.3.3）。在我们的类比下，这意味着 Plancherel 测度在 $\LS^{\dR,\mathrm{irred}}_{\bG}$ 上为常数，取值为 $\Vect$。在可约局部系统附近，存在与伴随丛的幂零水平截影相关的修正，它们在这里是 Arthur 的 $SL_2$ 的化身。

\pnum{}我们现在给出一个简单例子，以说明一些基本技术要点。

假设 $X = \mathbb{P}^1$ 且 $G = \mathbb{G}_m$。于是 $\Bun_{\mathbb{G}_m}$ 参数化 $\mathbb{P}^1$ 上的线丛，因此同构于 $\ZZ \times \mathbf{B}\mathbb{G}_m$：$\ZZ$-因子参数化线丛的次数，而 $\mathbf{B}\mathbb{G}_m$-因子编码了 $\mathbb{P}^1$ 上每个线丛都有自同构群 $\mathbb{G}_m$ 这一事实（在 $S$-族中以恰当方式理解）。

因此，
\[
  \Dmod(\Bun_{\mathbb{G}_m}(\mathbb{P}^1)) = \prod_{n \in \ZZ} \Dmod(\mathbf{B}\mathbb{G}_m).
\]
范畴 $\Dmod(\mathbf{B}\mathbb{G}_m)$ 可以相当显式地算出。令 $\pi: \Spec(k) \to \mathbf{B}\mathbb{G}_m$ 为结构映射，我们提醒它是一个光滑覆盖。函子 $\pi^!: \Dmod(\mathbf{B}\mathbb{G}_m) \to \Dmod(\Spec(k)) = \Vect$ 显然保守，并具有左伴随 $\pi_!$。由基变换，$\Vect$ 的 endofunctor $\pi^! \pi_!$ 由与 $C_{\dR}(\mathbb{G}_m)$（$\mathbb{G}_m$ 的 de Rham 上同调）作张量给出。此外，由 Barr--Beck 的一个简单形式，这个 endofunctor $\pi^! \pi_!$ 具有自然的单子结构，对应于 $C_{\dR}(\mathbb{G}_m)$ 上来自群结构 $\mathbb{G}_m$ 的代数结构；并且，诱导函子 $\Dmod(\mathbf{B}\mathbb{G}_m) \to C_{\dR}(\mathbb{G}_m)$-$\mathrm{mod}$ 是一个等价。

最后，当然，$C_{\dR}(\mathbb{G}_m)$ 是 $k$ 的一个平方零扩张，由一个位于上同调次数 $-1$ 的单个生成元 $\eta$ 扩张而成（别名：一个次数 $-1$ 生成元上的对称代数；一个生成元上的同调分次外代数）——这不过反映了圆周的同调在次数 $0$ 与 $1$ 上是 $1$ 维这一初等事实。

于是我们有：
\begin{equation}
  \Dmod(\Bun_{\mathbb{G}_m}(\mathbb{P}^1)) = \prod_{n \in \ZZ} (k \oplus k\eta)\text{-}\mathrm{mod}.
\end{equation}

天真地看，$\mathbb{P}^1$ 是单连通的，因此没有非平凡局部系统，所以人们可能期望 $\LS^{\dR}_{\bG}(\mathbb{P}^1)$ 等于 $\mathbf{B}G$（反映平凡局部系统的非平凡自同构群）。对 $\bG = \mathbb{G}_m$，我们将有 $\QCoh(\mathbf{B}\mathbb{G}_m) = \Rep(\mathbb{G}_m) = \prod_{n \in \ZZ} \Vect$，这与 (2.4) 很接近，但缺少次数 $-1$ 的生成元。

事实上，这是因为我们太天真了。\emph{导出}叠 $\LS^{\dR}_{\bG}(\mathbb{P}^1)$ 等于\footnote{这个公式例如来自把 $\mathbb{P}^1$ 上的局部系统看作两个标准开集 $\mathbb{A}^1 \subset \mathbb{P}^1$ 上的局部系统连同它们在交 $\mathbb{A}^1 \setminus 0$ 上的同构。注意 $\mathbb{A}^1$ 是可缩的，而不只是单连通的。} $\mathbf{B}\bG \times_{\bG} \mathbf{B}\bG$，对 $\bG = \mathbb{G}_m$ 就是 $\mathbf{B}\mathbb{G}_m \times (0 \times_{\mathbb{A}^1} 0)$。这里的纤维积必须在\emph{导出代数几何}意义下取。于是我们得到 $\QCoh(\LS^{\dR}_{\mathbb{G}_m}(\mathbb{P}^1)) = \prod_{n \in \ZZ} \QCoh(0 \times_{\mathbb{A}^1} 0)$。最后，我们注意 $0 \times_{\mathbb{A}^1} 0$ 是 $\Spec$ 的 $k \otimes_{k[t]} k$（张量积是导出的，即包含 $\mathrm{Tor}^{k[t]}_i(k,k)$ 这些群），这恰是我们先前看到的同一个平方零扩张 $C_{\dR}(\mathbb{G}_m)$。

我们注意到，底层的经典叠恢复了我们先前的朴素局部系统概念。

或者，人们可以如下看待导出代数几何的用处。对一般 $X$ 与 $\bG$，标准论证说 $\LS^{\dR}_{\bG}$ 在 $\bG$-局部系统 $\sigma$ 处的切空间是 $H^1_{\dR}(X, \check{\mathfrak{g}}_\sigma)$，即以伴随局部系统 $\sigma$ 为系数的第一个 de Rham 上同调。更一般地，我们应当期望切\emph{复形}是 $C_{\dR}(X, \check{\mathfrak{g}}_\sigma)[1]$。正如上面的例子所示，这个公式只有在把 $\LS^{\dR}_{\bG}$ 解释为\emph{导出}叠时一般才能成立。

总结：我们在几何朗兰兹的谱侧使用导出代数几何，\emph{因为它给出正确的答案}（不像经典代数几何），并且\emph{因为它产生模空间更易处理的无穷小几何}。

\pnumt{结论。}{} 上面我们简要讨论了算术朗兰兹，并对（de Rham）几何朗兰兹作了较长的介绍。

有一个关键区别，§2.1.1 已有所暗示：在算术朗兰兹中，我们研究\emph{原子对象}（不可约表示），而在\emph{几何}朗兰兹中，我们研究分子对象（自守函数空间的一个类比物）。人们可以把这种情形与圆周 $S^1$ 上的 Fourier 理论作比较：\emph{原子}理论说 $S^1$ 的（必然酉的）特征标与 $\ZZ$ 一一对应，但实际的 Fourier 理论说 $L^2(S^1)$ 是 $\ZZ$ 上 $1$ 维 Hilbert 空间（即 $L^2(\ZZ)$）的直接积分。在自守理论中，后者的类比物将是理想的，但我们所能触及的只是前者。

[2] 的一个出发点是试图解决这一差异，至少对函数域，至少在处处无分歧情形。最终，我们得到一个更接近几何朗兰兹猜想的算术视角。

我想强调：我们的工作并非唯一试图弥合这一鸿沟的工作；[32] 与 [11] 是密切相关的努力，而我们都受到 V.~Lafforgue 的突破 [22] 的启发。
